% ECONOMICS_PROBLEM: {"id": "C73-252", "title": "Learning in games and convergence", "jel_code": "C73", "source_id": "OP-0002"}
% !TeX program = xelatex
\documentclass[11pt,a4paper]{article}
\usepackage[margin=23mm]{geometry}
\usepackage{fontspec}
\setmainfont{Latin Modern Roman}
\usepackage{amsmath,amssymb,mathtools}
\usepackage{xcolor,enumitem,booktabs,longtable,array,xurl,hyperref}
\definecolor{muted}{HTML}{53636C}
\hypersetup{colorlinks=true,urlcolor=blue,linkcolor=blue}
\setlength{\parindent}{0pt}
\setlength{\parskip}{5pt}
\newcommand{\R}{\mathbb R}
\newcommand{\E}{\mathbb E}
\newcommand{\N}{\mathbb N}
\newcommand{\ind}{\mathbf 1}
\newcommand{\Var}{\operatorname{Var}}
\newcommand{\Cov}{\operatorname{Cov}}
\newcommand{\supp}{\operatorname{supp}}
\DeclareMathOperator*{\argmin}{arg\,min}
\DeclareMathOperator*{\argmax}{arg\,max}
\newcommand{\entrymeta}[3]{{\sffamily\small\color{muted}\textbf{\texttt{#1}}\quad|\quad #2\par #3\par}}
\newcommand{\Problemref}[1]{\texttt{#1}}
\newcommand{\JELref}[2]{\texttt{#2} (JEL \texttt{#1})}
\begin{document}
\section*{Learning in games and convergence}
\textbf{Problem ID:} \texttt{C73-252}\quad \textbf{JEL:} \texttt{C73}\par
% BEGIN ECONOMICS STATEMENT
\label{problem:OP-0002}
\label{atlas:prob:M2}

\noindent\textbf{Original atlas ID:} M2. \textbf{Formulation:} editorial specification of a research family.

\subsection*{Definitions and assumptions}

For a finite game $G$ with actions $A_i$ and bounded utilities $u_i\in[0,1]$, let $a^t=(a_i^t)_i$ be play at dates $t=1,2,\ldots$. Player $i$ uses an adapted randomized policy, depending only on a specified information filtration $\mathcal F_i^t$. Define empirical joint play $\mu_T=T^{-1}\sum_{t=1}^T\delta_{a^t}$. External regret and swap regret are respectively
\[
R_i^{\rm ext}(T)=\max_{b\in A_i}\sum_{t=1}^T[u_i(b,a_{-i}^t)-u_i(a^t)],\qquad
R_i^{\rm swap}(T)=\max_{\phi:A_i\to A_i}\sum_{t=1}^T[u_i(\phi(a_i^t),a_{-i}^t)-u_i(a^t)].
\]
Regret statements concern the positive part of these quantities. A coarse-correlated equilibrium has no profitable fixed-action deviation under its joint law; a correlated equilibrium has no profitable deviation conditional on an action recommendation.

\subsection*{Formal research question}

For a prespecified class $\mathcal L$ of heterogeneous learning policies, information structures, and permitted switching processes, identify the subclasses for which empirical play approaches the coarse-correlated or correlated-equilibrium set, and obtain valid finite-horizon rates. Separately ask whether the conditional mixed-action profile $p^t=(\Pr(a_i^t=\cdot\mid\mathcal F_i^t))_i$ approaches the Nash set under a clearly specified common history and conditional independence protocol. For example, seek $r_T\to0$ with
\[
\sup_{L\in\mathcal L}\mathbb E_L\!\left[\operatorname{dist}(\mu_T,\mathcal E(G))\right]\le r_T,
\]
where $\mathcal E(G)$ is the declared equilibrium set and distance uses a fixed norm. Determine when such a uniform statement fails and whether measurable switching restrictions restore it.

\subsection*{Interpretation and scope}

The distinction between empirical averages and last-iterate behavior is indispensable. Vanishing external regret implies that accumulation points of empirical play satisfy the coarse-correlated inequalities; vanishing swap regret supplies the correlated inequalities. Neither implication gives convergence of current play to Nash equilibrium. With finite actions, these statements follow by dividing the relevant deviation inequalities by $T$ and taking subsequential limits. Monitoring restrictions, changing payoffs, and switching between algorithms must enter $\mathcal L$; otherwise an arbitrary learner can cycle or keep choosing dominated actions. Empirical work should measure predictions at the same horizon and observation level as the mathematical result, and compare rules on held-out games.

\subsection*{Source audit and status}

Fudenberg--Levine (1998) is a verified monograph. Hart--Mas-Colell (2000) already proves almost-sure convergence of empirical play to the correlated-equilibrium set for its regret-matching procedure. Young (1993) studies a specific adaptive convention process. Thus convergence to correlated equilibrium is not itself an unsolved general benchmark, and these references do not imply last-iterate Nash convergence. The open component is a restricted characterization and empirically disciplined analysis of heterogeneous learning and switching. The class $\mathcal L$, uniform rate, and prediction protocol are editorial specifications; an actual theorem must state their restrictions and exhibit its proof or counterexample.

\subsection*{References}

\begin{enumerate}

\item[S1] Drew Fudenberg and David K. Levine (1998). \emph{The Theory of Learning in Games}. MIT Press. \url{https://mitpress.mit.edu/9780262061940/the-theory-of-learning-in-games/}. \textit{Locator:} publisher description; learning models. \textit{Establishes:} Theory relating adaptive behavior and equilibrium. \textit{Does not establish:} unconditional convergence of every adaptive process.

\item[S2] Sergiu Hart and Andreu Mas-Colell (2000). \emph{A Simple Adaptive Procedure Leading to Correlated Equilibrium}. Econometrica 68(5), 1127--1150. \url{https://onlinelibrary.wiley.com/doi/10.1111/1468-0262.00153}. \textit{Locator:} abstract; regret-matching convergence result. \textit{Establishes:} Almost-sure convergence of empirical play to the correlated-equilibrium set for the stated procedure. \textit{Does not establish:} last-iterate convergence to Nash equilibrium.

\item[S3] H. Peyton Young (1993). \emph{The Evolution of Conventions}. Econometrica 61(1), 57--84. \url{https://www.econ2.jhu.edu/people/young/Publications.html}. \textit{Locator:} article; perturbed adaptive coordination process. \textit{Establishes:} Stochastic selection of conventions under a particular learning process. \textit{Does not establish:} a universal behavioral law.

\end{enumerate}

\subsection*{Common conventions: Source mathematical conventions}

Unless an entry states otherwise, agent and firm sets are finite. State and action spaces are measurable, strategies and outcome maps are measurable, probability laws are specified, and expectations exist for the targets under discussion. Infinite-horizon discount factors lie in $(0,1)$ and discounted objectives are finite. Norms, tolerances and complexity input sizes must be specified for computational comparisons. Constants or primitives that appear locally are inputs, not empirical facts asserted by this catalogue.

\textbf{Identification.} A statistical model $m\in\mathcal M$ implies an observable law $P_m$ and target $\theta(m)$. At law $P$, its identified set is
\[
\Theta(P;\mathcal M)=\{\theta(m):m\in\mathcal M,\ P_m=P\}.
\]
Point identification holds when this set is a singleton, or equivalently when $P_{m_1}=P_{m_2}$ implies $\theta(m_1)=\theta(m_2)$ for all compatible models. Sharp bounds mean bounds attained, or approached when the boundary is not attained, by compatible models. Writing this set is a definition, not a solution: the substantive task is to establish informative restrictions and derive or estimate the set. An unrestricted-confounding model may give only trivial bounds.

\textbf{Inference.} Identification, estimability and valid uncertainty quantification are separate questions. Fix $\alpha\in(0,1)$. A confidence procedure $C_n$ over a declared law class $\mathcal P$ has asymptotic uniform coverage at level $1-\alpha$ when
\[
\liminf_{n\to\infty}\inf_{P\in\mathcal P}\Pr_P\{\theta(P)\in C_n\}\geq1-\alpha.
\]
For set-identified targets the coverage event must be defined separately, for example coverage of the full identified set. If an entry uses identification-indexed models instead of uniquely indexed laws, the same coverage convention is applied over those models. Pointwise limits do not establish uniform validity, and asymptotic validity does not guarantee finite-sample precision. Sample sizes, dependence structures and weak-identification sequences are part of each application.

\textbf{Counterfactuals.} $Y_i(a)$ is a potential outcome under a specified intervention, including its timing, implementation and versions. It is not $Y_i$ conditional on voluntarily choosing $a$. A full intervention vector permits interference; shorthand scalar treatment notation does not silently impose no interference. Assignment, selection, attrition, anticipation and measurement must be specified. A claim about an instrument's local effect is not automatically a population policy effect.

\textbf{Economic equilibrium.} A counterfactual model includes best responses, constraints, feasibility, market clearing, state transitions and expectations consistent with the stated information. When several equilibria exist, either a declared selector is an input or the target is set-valued. Comparisons across policy regimes must use the stated selection convention. Reduced-form relationships are not presumed invariant after an equilibrium-changing intervention.

\textbf{Optimization and welfare.} Optimization is over the stated feasible set. If attainment is not established, $\sup$ or $\inf$ and $\varepsilon$-optimal choices replace an asserted maximizer or minimizer. For example, with value $V=\sup_{a\in\mathcal A}W(a)<\infty$, an $\varepsilon$-optimal set is $\{a\in\mathcal A:W(a)\geq V-\varepsilon\}$ for $\varepsilon>0$. Benefits, costs, risk and health or environmental outcomes can be aggregated only using declared units and conversion weights. Interpersonal utility weights, the treatment of fiscal transfers and foreign profits, discounting, and the behavioral welfare criterion are explicit inputs. A comparative-static derivative additionally requires existence and appropriate differentiability of the selected counterfactual.

\textbf{Sources.} Local labels [S1], [S2], and so forth restart in every question. Locators identify the relevant abstract, model, section or result at the level actually checked; a bibliographic or abstract-level verification is not represented as inspection of an entire proof. Working papers are distinguished from later publications and changing versions. Source summaries are paraphrased. All URL checks and editorial formulations refer to the audit date stated above.
% END ECONOMICS STATEMENT
\end{document}
